\section{The knowledge bilattice operations}

\begin{proposition}[The FOUR bilattice operations]\label{thm:four-bilattice}\lean{Obsmetrologia.Charge.joinK}\leanok\uses{thm:four-charges, thm:predication, def:articulating}
The charge carrier has knowledge join and meet, truth join and meet, negation,
and the knowledge/truth comparisons. All are defined by Boolean slot algebra
on the same four-point carrier.
\end{proposition}

\begin{theorem}[Knowledge order is a lattice]\label{thm:joinK-lattice}\lean{Obsmetrologia.Charge.joinK_assoc}\leanok\uses{thm:four-bilattice}
Knowledge join and meet satisfy the bounded-lattice laws by finite exhaustion
over the four charges. The corresponding truth operations satisfy the same
lattice laws in the truth order.
\end{theorem}

\begin{theorem}[Distributivity and interlacing]\label{thm:bilattice-distrib}\lean{Obsmetrologia.Charge.meetK_joinK_distrib}\leanok\uses{thm:joinK-lattice, thm:valuation-meetK}
The FOUR bilattice is distributive, and truth/knowledge operations interlace:
a truth operation distributes across a knowledge operation on the same finite
carrier.
\end{theorem}

\begin{theorem}[The seeded bottom $\bot_k$]\label{thm:kBot-least}\lean{Obsmetrologia.Charge.kBot_leK}\leanok\uses{thm:four-bilattice}
Knowledge-bottom is the identity for knowledge join and the least element of
the knowledge order. This is the value assigned before evidence raises a cell.
\end{theorem}

\begin{theorem}[Negation links the two orders]\label{thm:negate-duality}\lean{Obsmetrologia.Charge.negate_leT}\leanok\uses{thm:four-bilattice}
Negation swaps the evidence slots. It preserves the knowledge order and
reverses the truth order, fixing the knowledge poles while exchanging the two
truth poles.
\end{theorem}
